\documentclass[12pt]{article}
\usepackage[utf8]{inputenc}
\usepackage[T1]{fontenc}
\usepackage{hyperref}
\title{Tupuri Language Resources for K-12 Curriculum Development}
\author{Compiled by Claude Sonnet 5 (claude-sonnet-5) for Andr\'as Kornai}
\date{}
\begin{document}
\maketitle
\section{Introduction}
Tupuri (ISO 639-3 tui) is an Atlantic-Congo language of northern Cameroon and southwestern Chad, written in the Latin alphabet. Speaker estimates in the results range from about 230,000 to 450,000, depending on the source.
\section{Primary Educational Resources}
No school textbooks or primers were found in the search results.
\section{Linguistic and Reference Materials}
\begin{itemize}
\item Wikipedia, Tupuri language (phonology, classification): \url{https://en.wikipedia.org/wiki/Tupuri_language}
\item Seignobos and Tourneux, \emph{Histoire des Toupouri} (French; history, not language teaching): \url{https://halshs.archives-ouvertes.fr/halshs-03269154/file/Seignobos_Tourneux_Histoire%20des%20Toupouri.pdf}
\end{itemize}
The results cite Ruelland's \emph{Dictionnaire Tupuri-Fran\c{c}ais-Anglais} (SELAF, 1988). No URL for it appeared, so it is only noted here.
\section{Religious and Community Resources}
Joshua Project pages (Cameroon people group and language): \url{https://www.joshuaproject.net/index.php/people_groups/15614/CM} and \url{https://www.joshuaproject.net/index.php/languages/tui}. Omniglot gives a Tupuri text of the Lord's Prayer: \url{https://omniglot.com/writing/tupuri.htm}.
\section{Digital and Audio Resources}
Omniglot mentions a Tupuri academy website. I did not open it and it is not listed as a resource.
\section{Curriculum Development Context}
The language is not known to be taught in schools. Documentation exists (a dictionary and a history), but I found no pedagogical materials, so a curriculum would need new work.
\end{document}